\section{轨迹问题}

\begin{example}
    已知抛物线 $C: y^2=2 p x$ ，斜率为 $\frac{2}{3}$ 的直线 $l$ 交抛物线于 $M, N$ 两点，且 $M(1,-2)$.
    \begin{enumerate}
        \item 求抛物线 $C$ 的方程；
        \item 若动点 $P$ 在抛物线 $C$ 上，$F$ 为拋物线的焦点，线段 $P F$ 的中点为 $Q$ ，求点 $Q$ 的轨迹方程；
        \item 试探究：抛物线 $C$ 上是否存在点 $P$ ，使得 $P M \perp P N$ ？若存在，求出 $P$ 点坐标；若不存在，请说明理由．
    \end{enumerate}
\end{example}

\newpage{}

\begin{example}
    点 $P$ 是椭圆 $\frac{x^2}{a^2} + \frac{y^2}{b^2} = 1 \quad (a>b>0)$ 上的一个动点，$F_1, F_2$ 是其两个焦点。若 $M$ 是 $\triangle PF_1F_2$ 的重心，则点 $M$ 的轨迹方程为 \underline{\hspace{4cm}}.
\end{example}

\begin{tips}
    回忆一下重心坐标公式。
\end{tips}

\v{12em}

\begin{example}
    设椭圆方程为 $\frac{x^2}{25} + \frac{y^2}{9} = 1$。$P$ 为椭圆上任意一点，$M$ 为 $x$ 轴上一点，且 $PM \perp x$ 轴。点 $N$ 在线段 $PM$ 上，且满足 $\overrightarrow{PN} = 2\overrightarrow{NM}$。求点 $N$ 的轨迹方程。
\end{example}

\v{12em}

\begin{problem}
已知在平面直角坐标系 $x O y$ 中，动点 $P$ 到 $(-\sqrt{3}, 0)$ 和 $(\sqrt{3}, 0)$ 的距离和为 4 ，设点 $A\left(1, \frac{1}{2}\right)$．
\begin{enumerate}
    \item 求动点 $P$ 的轨迹方程；
    \item $M$ 为线段 $P A$ 的中点，求点 $M$ 的轨迹方程；
    \item 过原点 $O$ 的直线交 $P$ 的轨迹于 $B, C$ 两点，求 $\triangle A B C$ 面积的最大值．
\end{enumerate}
\end{problem}

\v{12em}

\begin{problem}
已知椭圆 $E: \frac{x^2}{a^2}+\frac{y^2}{b^2}=1(a>b>0)$ 的左、右焦点分别为 $F_1, F_2$ ，左、右顶点分别为 $A, B, ~ P$ 是 $E$ 上异于 $A, B$ 的一个动点．若 $3\left|A F_1\right|=\left|B F_1\right|$ ，则下列说法正确的有
\begin{enumerate}
    \item 椭圆 $E$ 的离心率为 $\frac{1}{2}$
    \item 若 $P F_1 \perp F_1 F_2$ ，则 $\cos \angle P F_2 F_1=\frac{3}{5}$
    \item 直线 $P A$ 的斜率与直线 $P B$ 的斜率之积等于 $-\frac{3}{4}$
    \item 符合条件 $\overrightarrow{P F_1} \cdot \overrightarrow{P F_2}=0$ 的点 $P$ 有且仅有 2 个
\end{enumerate}
\end{problem}

\v{12em}

\begin{problem}
已知椭圆 $C: \frac{x^2}{a^2}+\frac{y^2}{b^2}=1(a>b>0)$ 的左、右两焦点分别是 $F_1$、$F_2$ ，其中 $\left|F_1 F_2\right|=2 c$ ．过左焦点的直线与椭圆交于 $A, B$ 两点．则下列说法中正确的有
\begin{enumerate}
    \item $\triangle A B F_2$ 的周长为 $4 a$
    \item 若 $A B$ 的中点为 $M, A B$ 所在直线斜率为 $k$ ，则 $k_{O M} \cdot k=-\frac{c^2}{a^2}$
    \item 若 $|A B|$ 的最小值为 $3 c$ ，则椭圆的离心率 $e=\frac{1}{3}$
    \item 若 $\overrightarrow{A F}_1 \cdot \overrightarrow{A F}_2=3 c^2$ ，则椭圆的离心率的取值范围是 $\left[\frac{\sqrt{5}}{5}, \frac{1}{2}\right]$
\end{enumerate}
\end{problem}

\v{12em}
